% Problem set: numbered problems with boxed final answers.
\documentclass[11pt,a4paper]{article}

\usepackage[T1]{fontenc}
\usepackage[margin=2.5cm]{geometry}
\usepackage{amsmath,amssymb,amsthm}
\usepackage{enumitem}
\usepackage{fancyhdr}

% Fill these in once; the header uses them on every page.
\newcommand{\coursename}{PHY 301: Quantum Mechanics}
\newcommand{\assignment}{Problem Set 1}
\newcommand{\studentname}{Your Name}

\pagestyle{fancy}
\fancyhf{}
\lhead{\coursename}
\chead{\assignment}
\rhead{\studentname}
\cfoot{\thepage}
\setlength{\headheight}{14pt}

% \problem{points} starts a numbered problem.
\newcounter{problem}
\newcommand{\problem}[1]{%
  \stepcounter{problem}%
  \section*{Problem \theproblem\ \normalfont\small(#1 points)}}
\newenvironment{solution}{\noindent\textbf{Solution.}\ }{\hfill$\square$\par\medskip}

\begin{document}

\begin{center}
  {\Large\bfseries \assignment} \\[4pt]
  \coursename \quad\textbullet\quad \studentname \quad\textbullet\quad Due \today
\end{center}

\problem{10}
Show that $\sum_{k=1}^{n} k = \dfrac{n(n+1)}{2}$.

\begin{solution}
Write the sum forwards and backwards and add the two lines term by term:
\begin{align*}
  S  &= 1 + 2 + \dots + n, \\
  S  &= n + (n-1) + \dots + 1, \\
  2S &= n(n+1).
\end{align*}
Hence
\[
  \boxed{S = \frac{n(n+1)}{2}}
\]
\end{solution}

\problem{15}
A particle is in the state $\psi(x) = A e^{-x^{2}/2a^{2}}$.
\begin{enumerate}[label=(\alph*)]
  \item Find the normalisation constant $A$.
  \item Compute $\langle x^{2} \rangle$.
\end{enumerate}

\begin{solution}
\begin{enumerate}[label=(\alph*)]
  \item Requiring $\int_{-\infty}^{\infty} |\psi|^{2}\,dx = 1$ gives
        $|A|^{2}\, a\sqrt{\pi} = 1$, so $A = (\pi a^{2})^{-1/4}$.
  \item $\langle x^{2} \rangle = a^{2}/2$.
\end{enumerate}
\end{solution}

\end{document}
